\documentclass[11pt,a4paper]{article}
\usepackage[utf8]{inputenc}
\usepackage[T1]{fontenc}
\usepackage[italian]{babel}
\usepackage{lmodern}
\usepackage{microtype}
\usepackage[margin=2cm,headheight=14pt]{geometry}
\usepackage{amsmath,amssymb,amsthm,mathtools,amsfonts,bm,siunitx}
\usepackage{booktabs,tabularx,enumitem,float}
\usepackage{xcolor}
\usepackage[most]{tcolorbox}
\usepackage{fancyhdr}
\usepackage{listings}
\usepackage{hyperref}
\usepackage{bookmark}

\definecolor{navyblue}{RGB}{10, 45, 90}
\definecolor{coolblue}{RGB}{28, 110, 164}
\definecolor{forestgreen}{RGB}{20, 115, 60}
\definecolor{warmamber}{RGB}{195, 110, 10}
\definecolor{crimsonred}{RGB}{175, 30, 45}
\definecolor{subtlegray}{RGB}{245, 247, 250}
\definecolor{codebg}{RGB}{240, 242, 245}

\hypersetup{
    colorlinks=true,
    linkcolor=navyblue,
    urlcolor=coolblue,
    citecolor=forestgreen
}

\pagestyle{fancy}
\fancyhf{}
\fancyhead[L]{\textbf{\textsf{\textcolor{navyblue}{Statistica I --- Soluzione Ufficiale Dettagliata}}}}
\fancyhead[R]{\textsf{\textcolor{coolblue}{III Appello (20/07/2022)}}}
\fancyfoot[C]{\thepage}
\renewcommand{\headrulewidth}{0.6pt}

\newtcolorbox{problemabox}[1]{
    enhanced,
    breakable,
    colback=white,
    colframe=navyblue,
    coltitle=white,
    fonttitle=\bfseries\sffamily,
    title={#1},
    attach boxed title to top left={yshift=-2mm, xshift=3mm},
    boxed title style={colback=navyblue, boxrule=0pt, sharp corners, rounded corners=north},
    boxrule=1.2pt,
    sharp corners,
    rounded corners=south,
    top=4mm, bottom=3mm, left=3mm, right=3mm
}

\newtcolorbox{soluzionebox}[1][Svolgimento Dettagliato]{
    enhanced,
    breakable,
    colback=subtlegray,
    colframe=forestgreen!80!black,
    coltitle=white,
    fonttitle=\bfseries\sffamily,
    title={#1},
    attach boxed title to top left={yshift=-2mm, xshift=3mm},
    boxed title style={colback=forestgreen!80!black, boxrule=0pt, sharp corners, rounded corners=north},
    boxrule=1pt,
    sharp corners,
    rounded corners=south,
    top=4mm, bottom=3mm, left=3mm, right=3mm
}

\newtcolorbox{esamebox}[1][Consiglio per l'Esame \& Aspetti Chiave]{
    enhanced,
    breakable,
    colback=crimsonred!4!white,
    colframe=crimsonred,
    coltitle=crimsonred,
    fonttitle=\bfseries\sffamily,
    title={#1},
    attach boxed title to top left={yshift=-1.5mm, xshift=2mm},
    boxed title style={colback=white, boxrule=0.8pt, colframe=crimsonred},
    boxrule=0.8pt,
    leftrule=3.5pt,
    sharp corners,
    top=3.5mm, bottom=2.5mm, left=3mm, right=3mm
}

\lstset{
    backgroundcolor=\color{codebg},
    basicstyle=\ttfamily\small,
    breaklines=true,
    frame=single,
    rulecolor=\color{navyblue!40},
    keywordstyle=\color{navyblue}\bfseries,
    commentstyle=\color{forestgreen}\itshape,
    stringstyle=\color{warmamber},
    showstringspaces=false
}

\newcommand{\EE}{\mathbb{E}}
\newcommand{\Prob}{\mathbb{P}}
\newcommand{\Var}{\operatorname{Var}}
\newcommand{\dx}{\mathrm{d}x}

\begin{document}

\begin{center}
    {\LARGE\bfseries\color{navyblue} Università di Pisa --- Corso di Laurea in Ingegneria Gestionale}\\[0.4em]
    {\Large\bfseries Statistica I (A.A. 2021/2022) --- III Appello del 20/07/2022}\\[0.3em]
    {\large\textsf{Soluzione Completa, Rigorosa e Commentata Passo-Passo}}
\end{center}
\vspace{0.5em}
\hrule height 1.2pt
\vspace{1.5em}

% ==============================================================================
% PROBLEMA 1
% ==============================================================================
\begin{problemabox}{Problema 1 (Testo Integrale)}
Consideriamo l'integrale
\[
\int_5^9 \frac{\log(2 + \cos(e^x))}{1 + x^3} \, \dx.
\]
\begin{enumerate}[label=(\alph*)]
    \item Discutere brevemente il metodo di integrazione di Monte Carlo nella sua forma generale e come può essere applicato per l'approssimazione di questo integrale in particolare.
    \item Scrivere i comandi in R la cui esecuzione restituisce un valore approssimativo dell'integrale.
\end{enumerate}
\end{problemabox}

\begin{soluzionebox}
\textbf{1. Metodo Monte Carlo su intervallo limitato $[a, b]$:}
Sia $I = \int_a^b g(x) \, \dx$. Moltiplicando e dividendo per la lunghezza dell'intervallo $(b - a)$:
\[
I = (b - a) \int_a^b g(x) \frac{1}{b - a} \, \dx = (b - a) \, \EE[g(X)], \qquad X \sim \operatorname{Unif}(a, b).
\]
Nel caso specifico, $a = 5$, $b = 9 \implies b - a = 4$, con $g(x) = \frac{\log(2 + \cos(e^x))}{1 + x^3}$.
Pertanto:
\[
I = 4 \, \EE\left[ \frac{\log(2 + \cos(e^X))}{1 + X^3} \right], \qquad X \sim \operatorname{Unif}(5, 9).
\]
Il valore numerico di riferimento è $I \approx \mathbf{0.008618}$.

\textbf{2. Implementazione in linguaggio R:}
\begin{lstlisting}[language=R]
set.seed(42)
N <- 1000000
x_samples <- runif(N, min = 5, max = 9)
g_val <- log(2 + cos(exp(x_samples))) / (1 + x_samples^3)
I_hat <- 4 * mean(g_val)
cat(sprintf("Stima Monte Carlo: %.6f\n", I_hat)) # ~ 0.0086
\end{lstlisting}
\end{soluzionebox}

% ==============================================================================
% PROBLEMA 2
% ==============================================================================
\begin{problemabox}{Problema 2 (Testo Integrale)}
Si consideri una variabile aleatoria $X$ che rappresenta la proporzione di elementi attivi in un processore, con densità:
\[
f_X(x) = \begin{cases}
c x^4 (1 - x^2), & x \in (0, 1), \\
0, & \text{altrimenti.}
\end{cases}
\]
Un processore è Functional (F) se $X \ge 0.7$, Recoverable (R) se $X \in (0.5, 0.7)$, Compromised (C) se $X \le 0.5$.
\begin{enumerate}[label=(\alph*)]
    \item Determinare il valore di $c$.
    \item Calcolare la probabilità che un processore sia Functional.
    \item Peschiamo 3 processori indipendenti: calcolare la probabilità di osservare i) la sequenza ordinata F, R, C e ii) almeno due processori Functional.
\end{enumerate}
\end{problemabox}

\begin{soluzionebox}
\begin{enumerate}[label=(\alph*)]
    \item \textbf{Normalizzazione:}
    \[
    1 = \int_0^1 c (x^4 - x^6) \, \dx = c \left[ \frac{x^5}{5} - \frac{x^7}{7} \right]_0^1 = c \left( \frac{1}{5} - \frac{1}{7} \right) = c \cdot \frac{2}{35} \implies \mathbf{c = \frac{35}{2} = 17.5}.
    \]

    \item \textbf{Probabilità dell'evento Functional ($p_F$):}
    \begin{align*}
    p_F = \Prob(X \ge 0.7) &= \frac{35}{2} \int_{0.7}^1 (x^4 - x^6) \, \dx = \frac{35}{2} \left[ \frac{x^5}{5} - \frac{x^7}{7} \right]_{0.7}^1 \\
    &= \frac{35}{2} \left( \frac{2}{35} - \left(\frac{0.7^5}{5} - \frac{0.7^7}{7}\right) \right) = 1 - \frac{35}{2}(0.033614 - 0.011765) \\
    &= 1 - \frac{35}{2}(0.021849) = 1 - 0.38236 = \mathbf{0.6176} \quad (\mathbf{61.76\%}).
    \end{align*}

    \item \textbf{Probabilità sui 3 processori estratti:}
    Calcoliamo preliminarmente le probabilità delle classi R e C:
    \begin{itemize}
        \item $p_C = \Prob(X \le 0.5) = \frac{35}{2} \left( \frac{0.5^5}{5} - \frac{0.5^7}{7} \right) = \frac{35}{2}(0.00625 - 0.001116) = \mathbf{0.0898}$.
        \item $p_R = \Prob(0.5 < X < 0.7) = 1 - p_F - p_C = 1 - 0.6176 - 0.0898 = \mathbf{0.2926}$.
    \end{itemize}
    \begin{enumerate}[label=\roman*)]
        \item \emph{Sequenza ordinata (F, R, C):}
        \[
        \Prob(F, R, C) = p_F \cdot p_R \cdot p_C = 0.6176 \cdot 0.2926 \cdot 0.0898 \approx \mathbf{0.0162} \quad (\mathbf{1.62\%}).
        \]
        \item \emph{Almeno due Functional su 3 (Binomiale $\operatorname{Bin}(3, p_F)$):}
        \[
        \Prob(\ge 2\text{ F}) = \binom{3}{2} p_F^2 (1 - p_F) + p_F^3 = 3(0.6176)^2(0.3824) + (0.6176)^3 = 0.4377 + 0.2356 = \mathbf{0.6733} \quad (\mathbf{67.33\%}).
        \]
    \end{enumerate}
\end{enumerate}
\end{soluzionebox}

% ==============================================================================
% PROBLEMA 3
% ==============================================================================
\begin{problemabox}{Problema 3 (Testo Integrale)}
Il numero di clienti giornalieri di un ristorante ha distribuzione di Poisson con parametro $\lambda_1 = 10$ nei giorni feriali (Mar--Gio, 3 giorni) e parametro $\lambda_2 = 16$ nel weekend (Ven--Dom, 3 giorni). Il ristorante è chiuso il lunedì.
\begin{enumerate}[label=(\alph*)]
    \item Se sappiamo che un certo giorno sono arrivati 7 clienti, qual è la probabilità che fosse un giorno del weekend?
    \item Identificare la distribuzione del numero di clienti in una settimana.
    \item Approssimare la probabilità che il ristorante abbia almeno 4000 clienti in un anno (52 settimane).
\end{enumerate}
\end{problemabox}

\begin{soluzionebox}
\begin{enumerate}[label=(\alph*)]
    \item \textbf{Aggiornamento Bayesiano per il giorno di arrivo:}
    I giorni di apertura sono 6 (3 feriali, 3 weekend), con prior equiprobabile: $\Prob(\text{Feriale}) = 3/6 = 1/2, \quad \Prob(\text{Weekend}) = 3/6 = 1/2$.
    Le verosimiglianze di osservare 7 clienti sono:
    \[
    \Prob(X=7 \mid \text{Feriale}) = \frac{e^{-10} 10^7}{7!}, \qquad \Prob(X=7 \mid \text{Weekend}) = \frac{e^{-16} 16^7}{7!}.
    \]
    Applicando il Teorema di Bayes:
    \begin{align*}
    \Prob(\text{Weekend} \mid X=7) &= \frac{\frac{e^{-16} 16^7}{7!} \cdot \frac{1}{2}}{\frac{e^{-16} 16^7}{7!} \cdot \frac{1}{2} + \frac{e^{-10} 10^7}{7!} \cdot \frac{1}{2}} = \frac{e^{-16} 16^7}{e^{-16} 16^7 + e^{-10} 10^7} \\
    &= \frac{1}{1 + e^6 \left(\frac{10}{16}\right)^7} = \frac{1}{1 + 403.4288 \cdot 0.03725} \approx \mathbf{0.0624} \quad (\mathbf{6.24\%}).
    \end{align*}

    \item \textbf{Distribuzione settimanale dei clienti:}
    La somma di variabili di Poisson indipendenti è ancora una variabile di Poisson con parametro pari alla somma dei parametri:
    \[
    \lambda_{\text{sett}} = 3 \lambda_1 + 3 \lambda_2 = 3(10) + 3(16) = 30 + 48 = \mathbf{78}.
    \]
    Il numero di clienti settimanali $W \sim \mathbf{\operatorname{Poiss}(\lambda = 78)}$.

    \item \textbf{Clienti annuali e approssimazione con TLC ($N = 52$ settimane):}
    Il numero totale annuo di clienti $V$ segue una legge $\operatorname{Poiss}(52 \times 78 = 4056)$:
    \[
    \EE[V] = 4056, \qquad \Var(V) = 4056 \implies \sigma_V = \sqrt{4056} \approx 63.6867.
    \]
    Per il TLC, $V \approx \mathcal{N}(4056, 4056)$. Standardizzando:
    \[
    \Prob(V \ge 4000) \approx 1 - \Phi\left(\frac{4000 - 4056}{\sqrt{4056}}\right) = 1 - \Phi\left(\frac{-56}{63.6867}\right) = 1 - \Phi(-0.8793) = \Phi(0.8793).
    \]
    Dalle tavole della normale standard ($\Phi(0.88) = 0.81057$):
    \[
    \Prob(V \ge 4000) \approx \mathbf{0.8104} \quad (\mathbf{81.04\%}).
    \]
\end{enumerate}
\end{soluzionebox}

% ==============================================================================
% PROBLEMA 4
% ==============================================================================
\begin{problemabox}{Problema 4 (Testo Integrale)}
La media storica dell'escursione termica a Roma è $13.0^\circ\text{C}$. Nel 2021 vengono effettuate 8 misurazioni:
\[
13.7, \quad 13.5, \quad 13.3, \quad 12.9, \quad 13.0, \quad 13.4, \quad 13.2, \quad 12.8.
\]
\begin{enumerate}[label=(\alph*)]
    \item Si calcoli l'intervallo di confidenza bilaterale al 95\% per $\mu$.
    \item Si formulino le ipotesi $H_0$ e $H_1$ per mostrare che l'escursione termica media è aumentata.
    \item Su 365 misurazioni si ottengono $\overline{x} = 13.08$ e $s^2 = (0.5)^2$: testare l'ipotesi a $\alpha = 2\%$.
    \item Quante misurazioni sarebbero necessarie per confermare il claim a $\alpha = 1\%$?
\end{enumerate}
\end{problemabox}

\begin{soluzionebox}
\textbf{1. Statistiche descrittive ($n = 8$):}
\[
\overline{x} = \mathbf{13.225^\circ\text{C}}, \qquad s = \mathbf{0.3105^\circ\text{C}}.
\]

\textbf{2. Risoluzione Quesito (a) --- Intervallo di Confidenza al 95\% ($\nu = 7$ g.d.l.):}
Il valore critico è $t_{7, \, 0.025} = 2.3646$.
\[
ME = 2.3646 \cdot \frac{0.3105}{\sqrt{8}} = 2.3646 \cdot 0.10978 = \mathbf{0.2596^\circ\text{C}}.
\]
\[
\mathbf{\operatorname{IC}_{95\%}(\mu) = [13.225 - 0.2596, \;\; 13.225 + 0.2596] = [12.9654, \;\; 13.4846]^\circ\text{C}}.
\]

\textbf{3. Risoluzione Quesito (b) --- Ipotesi statistiche:}
\[
\mathbf{H_0: \mu \le 13.0^\circ\text{C}} \qquad \text{vs} \qquad \mathbf{H_1: \mu > 13.0^\circ\text{C}}.
\]

\textbf{4. Risoluzione Quesito (c) --- Test a grandi campioni ($n = 365, \alpha = 0.02$):}
\[
Z_{\text{oss}} = \frac{\overline{x} - \mu_0}{s/\sqrt{n}} = \frac{13.08 - 13.00}{0.50/\sqrt{365}} = \frac{0.08}{0.02617} = \mathbf{3.0568}.
\]
Il valore critico unilaterale destro a $\alpha = 0.02$ è $z_{0.02} = 2.0537$.
Poiché $Z_{\text{oss}} = 3.0568 > 2.0537$, \textbf{rifiutiamo l'ipotesi nulla $H_0$} ($p$-value $= 1 - \Phi(3.0568) = \mathbf{0.0011} < 2\%$).
\emph{Conclusione:} I dati smentiscono fortemente chi nega il cambiamento climatico, fornendo evidenza rigorosa di un aumento della temperatura media.

\textbf{5. Risoluzione Quesito (d) --- Numerosità campionaria necessaria per $\alpha = 1\%$:}
Il quantile critico a $\alpha = 0.01$ è $z_{0.01} = 2.3263$.
\[
\frac{13.08 - 13.00}{0.50 / \sqrt{n}} \ge 2.3263 \implies \sqrt{n} \ge \frac{2.3263 \cdot 0.50}{0.08} = 14.5394 \implies n \ge (14.5394)^2 \approx 211.39.
\]
Sono necessarie almeno $\mathbf{n = 212}$ misurazioni.
\end{soluzionebox}

% ==============================================================================
% PROBLEMA 5
% ==============================================================================
\begin{problemabox}{Problema 5 (Testo Integrale)}
Sia dato un campione $X_1,\dots,X_n$ con funzione di massa
\[
\Prob_\theta(X = k) = \frac{1}{\theta}\left(1 - \frac{1}{\theta}\right)^{k-1}, \qquad k \in \{1, 2, 3, \dots\}, \quad \theta \in (1, \infty).
\]
\begin{enumerate}[label=(\alph*)]
    \item Verificare che sia una distribuzione di probabilità per ogni $\theta > 1$. Di quale legge si tratta?
    \item Calcolare lo stimatore di massima verosimiglianza $\widehat{\theta}$.
    \item Esiste $n$ tale che questo stimatore non è corretto?
\end{enumerate}
\end{problemabox}

\begin{soluzionebox}
\begin{enumerate}[label=(\alph*)]
    \item \textbf{Identificazione e Normalizzazione:}
    Posto $p = 1/\theta \in (0, 1)$, riconosciamo la legge \textbf{Geometrica} supportata su $\mathbb{N}_{>0}$:
    \[
    \sum_{k=1}^\infty \Prob_\theta(X=k) = \frac{1}{\theta} \sum_{k=1}^\infty \left(1 - \frac{1}{\theta}\right)^{k-1} = \frac{1}{\theta} \cdot \frac{1}{1 - (1 - 1/\theta)} = \frac{1}{\theta} \cdot \theta = \mathbf{1}.
    \]

    \item \textbf{Stimatore di Massima Verosimiglianza (MLE):}
    La funzione di verosimiglianza su $k_1, \dots, k_n$ è:
    \[
    L(\theta) = \prod_{i=1}^n \left[ \frac{1}{\theta}\left(1 - \frac{1}{\theta}\right)^{k_i - 1} \right] = \theta^{-n} \left(1 - \frac{1}{\theta}\right)^{\sum_{i=1}^n k_i - n}.
    \]
    Prendendo il logaritmo:
    \[
    \ell(\theta) = -n \log \theta + \left(\sum_{i=1}^n k_i - n\right) \log\left(1 - \frac{1}{\theta}\right).
    \]
    Derivando rispetto a $\theta$:
    \[
    \frac{\partial \ell}{\partial \theta} = -\frac{n}{\theta} + \left(\sum_{i=1}^n k_i - n\right) \frac{1/\theta^2}{1 - 1/\theta} = -\frac{n}{\theta} + \frac{\sum_{i=1}^n k_i - n}{\theta(\theta - 1)} = 0.
    \]
    Risolvendo per $\theta$:
    \[
    n(\theta - 1) = \sum_{i=1}^n k_i - n \implies n\theta - n = \sum_{i=1}^n k_i - n \implies \mathbf{\widehat{\theta}_{\text{MLE}} = \frac{1}{n}\sum_{i=1}^n X_i = \overline{X}_n}.
    \]

    \item \textbf{Verifica della correttezza:}
    Per una variabile geometrica con parametro $p = 1/\theta$, il valore atteso è $\EE[X] = 1/p = \theta$.
    \[
    \EE[\widehat{\theta}_{\text{MLE}}] = \EE[\overline{X}_n] = \frac{1}{n}\sum_{i=1}^n \EE[X_i] = \theta.
    \]
    Poiché $\EE[\widehat{\theta}] = \theta$ per ogni $\theta > 1$ e per ogni dimensione campionaria $n \ge 1$, lo stimatore è \textbf{sempre corretto} ($\mathbf{\operatorname{Bias} = 0 \quad \forall n}$). Non esiste alcun valore di $n$ per cui sia distorto.
\end{enumerate}
\end{soluzionebox}

\end{document}
